\unnumberedChapter{Processor Abstraction Layer}
This chapter presents the Twin-64 processor abstraction layer....

\unnumberedSection{Concepts}

\unnumberedSection{Design Rationale}

\unnumberedSection{PAL Functions}




PAL\_TOC - programatic reset

PAL\_CHECK - entry point to HPMC

PAL\_RESET - hardware entry point, start of PAL code area.

PAL\_CHASSIS - writes a value to system chassis display.

PAL\_GETC - reads a character from console

PAL\_PUTC - writes a character to console

PAL\_GET\_LINE - reads a line from console

PAL\_PUT\_LINE - writes a line to console

PAL\_CONSOLE - general entry point to console functions. ( getc, putc, flus, write line, read line, etc. )


\unnumberedSection{PAL code area layout}

first locations are RESET vector, PAL COMMAND vector, two reserved vectors.

RESET vector is where a processor starts after reset.

COMMAND vector is the entry point to PAL functions.



\unnumberedSection{System startup}

PAL plays a key role in startup.


\begin{smallGapItemize}

\item{Test processor}

\item{Select monarch. All other processors wait for monarch. The monarch is the processor with the lowest HPA address. }

\item{Monarch issues RESET to all I/O modules.}

\item{Monarch finds memory module with largest size. If there are more than one, the lowest HPA wins.}

\item{Monarch places this module at physical memory address zero.}

\item{Monarch checks source of RESET. ( power fail ?) }

\item{Monarch tests memory large enough for ISL image }

\item{Monarch finds CONSOLE module }

\item{Monarch finds boot device }

\item{Monarch loads and launches IPL }


	
\end{smallGapItemize}

Page Zero plays a key role, lots of low level data, such as boot paths, etc.
Page Zero also has initial stack ? 

Destructive vs. non-destructive memory test ?


